\documentclass[sigconf,anonymous,review]{acmart}
% FSE 2027 — 10 pages main text + 2 pages references. Double-anonymous.

\begin{document}

\title{Prompt Engineering Cannot Fake Memorization: A Counterfactual Framework for
Measuring Exposure-Specific Code Reproduction in LLMs}

\author{Anonymous Author(s)}
\affiliation{\institution{Anonymous Institution}\country{}}

\begin{abstract}
Current evaluations of code memorization compare model outputs against target code
without controlling for what an unexposed model would produce under the same prompt,
conflating \emph{convergence} (an informative prompt narrows the solution space) with
\emph{recall} (the model retrieves target-specific information seen in training). We
introduce a controlled counterfactual framework whose primary outcome is the
\emph{exposed--unexposed gap}: the difference in performance between fine-tuning targets
and structurally matched targets held out of fine-tuning. Using post-cutoff code as a
clean exposure-free baseline, deliberate fine-tuning at frequencies
$k\in\{1,5,25\}$, and a matched unexposed holdout, we evaluate a ten-strategy prompt
taxonomy across code reproduction, source attribution, and high-entropy canary recovery.
We identify which prompting strategies increase the exposure-specific signal rather than
merely raising raw similarity.
\end{abstract}

\maketitle

\section{Introduction}                 % ~0.7 pp — report §16
% convergence problem w/ concrete example; gap; three contributions

\section{Background and Related Work}   % ~0.8 pp
% Carlini 2019/21/23; Kandpal 2022; Zhang 2023 (closest competitor); Al-Kaswan 2024;
% Lopes 2017 / Lee 2022 duplication; Meeus 2024 canaries; Karamolegkou 2023 license

\section{Framework}                     % ~0.8 pp
% T_E / T_U / T_H; Delta_pi; dose-response; three tasks and their SNR properties

\section{Experimental Design}           % ~1.2 pp
% corpus construction (post-cutoff source, membership + timestamp oracles, dedup,
% matching, canary injection); models + LoRA protocol; prompt taxonomy table;
% metrics; analysis plan

\section{Results: Exposure, Convergence, Interaction (RQ1--RQ3)}  % ~2.0 pp
% Delta_pi heatmap (primary model); convergence figure M(T_U, pi) base model;
% interaction table Delta_pi - Delta_{P1a}

\section{Results: Task Dependence and Dose-Response (RQ4--RQ5)}   % ~1.5 pp
% cross-task best strategies; dose-response curves; canaries cleanest, license weakest

\section{Results: Pre-cutoff and License (RQ6--RQ7)}             % ~0.8 pp

\section{Discussion}                    % ~0.8 pp
% copyright red-teaming implications; raw similarity is not memorization; fine-tuning proxy

\section{Threats to Validity}           % ~0.5 pp — report §15
% construct: fine-tuning vs pretraining; internal: clone confound, H2, test contamination;
% external: post-cutoff proxy, model families; conclusion: multiple comparisons, small k

\section{Data Availability}
% see DATA_AVAILABILITY.md

\section{Conclusion}                    % ~0.2 pp

\bibliographystyle{ACM-Reference-Format}
\bibliography{references}

\end{document}
